\chapter{Literature Review}
\label{chap:lit}

This chapter situates EdgeVerify within three bodies of work: (i)
self-supervised representation learning and the problem of representational
collapse, (ii) joint-embedding predictive architectures, and (iii) edge,
Dew, and decentralized computing. For each, we identify the mechanisms
relevant to the proposed design and the gaps that motivate it.

\section{Self-Supervised Representation Learning}
Self-supervised learning (SSL) removes the dependence on labeled data by
constructing a pretext task from the input itself. Contrastive methods
such as SimCLR \cite{chen2020simple} learn invariances by pulling together
embeddings of augmented views of the same image while pushing apart
embeddings of different images. Their principal limitation is a reliance
on large batches of negative samples to avoid trivial solutions---a
requirement that is poorly suited to memory-constrained edge hardware.

A subsequent line of work removed negatives entirely. BYOL
\cite{grill2020bootstrap} trains an online network to predict the output of
a momentum-updated target network, demonstrating that a predictor combined
with a stop-gradient and an Exponential Moving Average (EMA) target can
prevent collapse without contrastive pairs. SimSiam \cite{chen2021exploring}
further showed that the stop-gradient operation, rather than the momentum
encoder itself, is the critical component preventing collapse. EdgeVerify
adopts this EMA-target-plus-predictor design (Chapter~\ref{chap:method}),
but pairs it with explicit statistical regularization rather than relying
on architectural asymmetry alone.

\section{Regularization Against Dimensional Collapse}
An alternative to asymmetric, negative-free designs is to regularize the
statistics of the embedding space directly. Barlow Twins
\cite{zbontar2021barlow} minimizes redundancy by driving the
cross-correlation matrix between two views toward the identity. VICReg
\cite{bardes2022vicreg} decomposes the objective into three explicit
terms: an \emph{invariance} term (mean-squared error between views), a
\emph{variance} term that maintains the standard deviation of each
embedding dimension above a hinge threshold, and a \emph{covariance} term
that decorrelates dimensions. This explicit, non-contrastive formulation
is directly applicable at the level of individual layers and small
batches, which is why EdgeVerify's loss adapts the variance and covariance
terms of VICReg. The distinction in this work is that the regularizers are
applied within a \emph{local, layer-wise distillation} objective rather
than a single global loss.

\section{Analyses of Representation Collapse and Decorrelation}
\label{sec:collapse-lit}
A parallel body of work analyses collapse itself. Jing et al.
\cite{jing2022understanding} give a linear-model account of
\emph{dimensional collapse} in contrastive SSL, showing that embeddings can
occupy a low-dimensional subspace even when a global collapse is avoided.
Hua et al. \cite{hua2021feature} identify feature decorrelation as the key
mechanism and propose explicit decorrelation to counteract it. Most directly
relevant to this thesis, two works study \emph{where} decorrelation must act
in non-contrastive (BYOL/SimSiam-style) architectures: Li et al.
\cite{li2022understanding} and Liu et al. \cite{liu2022bridging} connect the
predictor and stop-gradient asymmetry to an implicit decorrelation and
analyse the conditions under which the representation collapses.

This thesis is closest to that last group, and it is important to state the
relationship precisely: prior work has \emph{already established}, for
Siamese non-contrastive methods, that the placement of the decorrelation
mechanism governs collapse. The present work does not discover that
phenomenon; it extends it along three axes that the prior analyses do not
cover. First, the setting is a \emph{predictive} JEPA with an
action-conditioned predictor and an \emph{explicit} VICReg-style
variance-plus-covariance objective, rather than a Siamese network with
implicit decorrelation. Second, we show that the failure is \emph{invisible}
to both the training loss and the VICReg variance criterion---a diagnostic
observation not emphasised in that work---and validate representation quality
with a downstream linear probe. Third, we test whether the correct placement
\emph{transfers to decentralized (federated) training}. The contribution is
therefore an extension of an established collapse phenomenon to a new
architecture and training regime, not a new phenomenon.

\section{Joint-Embedding Predictive Architectures}
LeCun's position paper on autonomous machine intelligence
\cite{lecun2022path} argues for prediction in representation space as a
route to efficient world models, framing the JEPA as a core building
block. I-JEPA \cite{assran2023self} instantiated this for images:
given a context block, a predictor estimates the representations of
masked target blocks produced by a target encoder, with no pixel-level
reconstruction. This yields strong representations while avoiding the
computational cost of generative decoding, and---unlike reconstructive
masked autoencoders \cite{he2022masked}---never materializes the
high-dimensional output. EdgeVerify inherits this predict-in-latent-space
principle and extends it with an \emph{action-conditioned} predictor:
the prediction is conditioned on a spatial action (a bounding box),
aligning the architecture with the conditional structure of world models
in \cite{lecun2022path}.

\section{Knowledge Distillation}
The teacher--student framing of EdgeVerify is related to knowledge
distillation \cite{hinton2015distilling}, in which a student network is
trained to match the outputs of a teacher. In the SSL setting, DINO
\cite{caron2021emerging} recast self-distillation with no labels as an
effective pretraining strategy for vision transformers, using a
momentum teacher and centering/sharpening to avoid collapse. EdgeVerify's
EMA target encoder plays the role of a slowly-evolving teacher, but the
anti-collapse mechanism is the VICReg-style variance/covariance penalty
rather than centering and sharpening.

\section{Edge, Dew, and Decentralized Computing}
Edge computing relocates computation from centralized clouds toward the
data source to reduce latency and bandwidth and to improve privacy
\cite{shi2016edge}. Dew computing pushes this further, positioning
independent, on-device capability that can function with minimal or
intermittent connectivity to higher tiers \cite{wang2016dew}. Decentralized
learning approaches such as federated learning \cite{mcmahan2017communication}
keep raw data local and communicate only model updates, mitigating privacy
risk at the cost of communication rounds. EdgeVerify targets a stricter
regime than federated learning: it aims to keep both data \emph{and}
training local to a single resource-constrained node, which sharpens the
memory-budget problem and motivates the layer-wise gradient isolation
described in the next chapter.

\section{Summary and Research Gap}
The reviewed literature establishes (i) that negative-free SSL is feasible
via EMA targets and predictors, (ii) that explicit variance/covariance
regularization prevents collapse without negatives, (iii) that JEPAs
avoid the cost of generative decoding, and (iv) that, in Siamese
non-contrastive methods, the placement of the decorrelation mechanism
governs collapse \cite{li2022understanding, liu2022bridging}. What remains
underexplored---and what this thesis addresses---is (a) whether the same
placement effect holds in an \emph{explicitly} regularized predictive JEPA,
(b) that the failure is invisible to the training loss and the variance
criterion and must be detected by effective rank and a downstream probe, and
(c) whether the corrected placement survives \emph{decentralized} training on
edge nodes. The thesis therefore extends an established collapse phenomenon
to the JEPA and federated settings, using an action-conditioned architecture
intended for on-device execution.
